\chapter{Literature Search Records and Evidence Tables}
\label{app:literature}
\section{Discovery Population and Evidence Status}
The screening record dated 16 November 2026 lists 120 unique candidates. Of these, 39 met the inclusion criteria, 37 were excluded and 44 remained potentially relevant but insufficiently verified. The corresponding full-text appraisal covered 37 of the 39 included works. The two accessible only through abstracts and bibliographic metadata contributed to the screening counts but supplied no methodological or numerical findings.\ev{Historical search and screening record}{\EvidenceRef{E006}{}, \EvidenceRef{E009}{}, \EvidenceRef{E005}{}}

The subsequent appraisal covered 59 PDF files. Two pairs were byte-identical, leaving \ExpandedPDFCount{} distinct files. Titles, authors, DOIs and repository identifiers showed that these represented \ExpandedPDFCount{} additional works, bringing the total appraised to \ReviewWorkCount{}. This historical total includes empirical industrial studies, transferable methods, statistical and representation methods, dataset descriptions and contextual studies. The 25 November revision adds one separately identified methodological reference, \citet{Ting1999Issues}, to address the missing original foundation for stacked generalization. It is stored in \texttt{references\_additions.bib} and is outside the two-stage appraisal counts and matrices; the review was not restarted.

The search log for this screening round records more than 15 distinct queries. Several queries returned results together, precluding an exact count of the records retrieved by each. The subsequent appraisal drew on additional accessible full texts and used publisher or repository pages to check bibliographic identities and versions. This procedure limits the review's coverage to the works identified and accessible through these routes.\ev{Expanded appraisal and counts}{\EvidenceRef{E016}{}, \EvidenceRef{E010}{}, \EvidenceRef{E011}{}, \EvidenceRef{E012}{}}

\section{Extraction, Locators, and Versions}
For each study, the appraisal recorded the prediction unit and target, physical inputs, derived representations, fusion method, single-input comparisons and assignment of data to training and evaluation. Missing inputs, noisy inputs and changes of operating conditions were coded separately. Numerical findings cite the relevant page, table or figure. Where the available text left a procedure unspecified, its status was recorded as unknown.

The appraisal distinguishes the published bibliographic record from the version examined. Preprints and journal articles sometimes give different sample counts for the same work. File hashes identify the version used for each extraction. The 37 initial full-text records include 148 short quotations that locate the evidence for their codes. The review compared reported methods and results without rerunning the published experiments.\ev{Extraction and identity records}{\EvidenceRef{E007}{}, \EvidenceRef{E008}{}, \EvidenceRef{E017}{}, \EvidenceRef{E014}{}}

\section{Comparison Matrix and Denominators}
Table~\ref{tab:literature-matrix} gives the codes for the initial 37 studies. Table~\ref{tab:expanded-register} identifies the contribution and principal limitation of each of the 57 additional works. The accompanying matrix gives the missing-input, noise and shift tests for individual studies, together with source locations.\ev{Full expanded comparison register}{\EvidenceRef{E013}{}, \EvidenceRef{E015}{}}

Within-protocol input variants do not guarantee equal capacity, pretraining, calibration or selection budget. Planned single-input deployment, missing training images, generated substitution, stochastic dropout and unexpected inference outage are different conditions. Ordinary defective objects, natural geometric gaps and fixed noisy data are not automatically controlled corruption experiments. The codes describe this inspected evidence and cannot estimate field-wide prevalence.

\begin{longtable}{@{}p{0.33\textwidth}llllll@{}}
\caption{Comparison of the 37 Original Full-Text Studies}\label{tab:literature-matrix}\\
\toprule Study & Data & Fusion & Base & Missing & Noise & Shift \\
\midrule\endfirsthead
\toprule Study & Data & Fusion & Base & Missing & Noise & Shift \\
\midrule\endhead
\citep{Wang2025A} & VIB & H & U & N & S & A \\
\citep{Ying2025Trustworthy} & CAM-S & H & M & N & S & N \\
\citep{Huang2026Rolling} & AUD-VIB & I & M & N & S & C \\
\citep{Gu2024Improved} & AUD-VIB & I & P & N & N & N \\
\citep{Yang2024Fault} & MS & I & M & PA & N & N \\
\citep{Ma2025A} & IMG-U & I & M & TR & N & N \\
\citep{Li2025PGMTKD} & MS & KD & R & PA & N & N \\
\citep{AlDulaimi2019A} & SV & I & L & N & S & N \\
\citep{Li2025MADFlow} & RGB-G & I & A & N & N & N \\
\citep{Wan2024Pipeline} & AUD & E & M & N & N & N \\
\citep{Baltrusaitis2019Multimodal} & X & TAX & S & NA & NA & NA \\
\citep{Sui2026Incomplete} & RGB-G & KD+L & M & PA & N & N \\
\citep{Cheng2025Multimodal} & RGB-G & H & M & N & N & N \\
\citep{Xu2025ModalityResilient} & RGB-G & KD+H & P & SD & N & N \\
\citep{Shangguan2025CPIR} & RGB-G & H & A & N & N & N \\
\citep{Altinses2025Performance} & SIM-CS & I & A & GS & N & T \\
\citep{Wang2023Multimodal} & RGB-G & H & M & N & N & N \\
\citep{Horwitz2023Back} & RGB-G & E & M & N & N & N \\
\citep{Costanzino2024Multimodal} & RGB-G & H & A & N & N & N \\
\citep{Rudolph2023Asymmetric} & RGB-G & I & M & N & N & N \\
\citep{Bonfiglioli2023The} & SIM-G & E & M & N & N & N \\
\citep{Bergmann2022The} & RGB-G & E & P & N & N & N \\
\citep{Wang2020What} & X & I & M & N & N & N \\
\citep{Xie2025Audiovisual} & CAM-S & I+KD & M & PA & S & N \\
\citep{Goodarzi2023Deep} & MS & E+I & M & N & S & N \\
\citep{Roth2022Towards} & RGB & U & U & N & N & N \\
\citep{Bergmann2019MVTec} & RGB & U & U & N & N & N \\
\citep{Purohit2019MIMII} & AUD & U & U & N & B & N \\
\citep{Koizumi2019ToyADMOS} & AUD & U & U & N & B & N \\
\citep{Koizumi2020Description} & AUD & U & U & N & F & N \\
\citep{Chu2024A} & AUD-VIB & I & N & N & N & A \\
\citep{Kullu2022A} & MS & I & M & N & N & N \\
\citep{Chen2023Multisensor} & CAM-S & E & M & N & N & N \\
\citep{Li2022A} & OPT-AUD & I & M & N & N & N \\
\citep{Petrich2021Multimodal} & CAM-S & I & M & RS & N & N \\
\citep{Pandiyan2022Deep} & OPT-AUD & E & N & N & N & W \\
\citep{Bevans2023Heterogeneous} & MULTI-CAM & E & A & N & N & B \\
\bottomrule
\end{longtable}
\displaynote{All 37 records have evidence status F. The codes describe the inspected protocol. They do not rate study quality.}
\noindent\textit{Data:} CAM-S denotes a physical camera and a separate non-image signal; RGB-G denotes RGB and geometry; OPT-AUD denotes photodiode and acoustic signals, without a camera input. MS denotes machine/process channels; SV denotes parallel representations of the same telemetry; IMG-U has incompletely documented image provenance. SIM-CS/SIM-G are synthetic. RGB, AUD and VIB are single physical modalities; MULTI-CAM uses several imaging sources; X is transferable non-industrial evidence.
\par\noindent\textit{Fusion:} E, early; I, intermediate; L, output-score fusion; H, hybrid feature/output fusion; KD, privileged cross-modal training; U, unimodal; TAX, taxonomy. Authors' terminology is normalized by where information is combined.
\par\noindent\textit{Baseline:} M, input variants evaluated within one protocol; P, partial/ambiguous; R, retained-input student with/without teachers; L, literature comparison; A, module/map/task ablations only; U, unimodal reference; S, survey; N, not established. M does not guarantee matched capacity, pretraining or tuning.
\par\noindent\textit{Missing:} PA, planned missing-input inference; SD, simulated dropout; TR, missing training images with complete tests; GS, generated substitution; RS, retrained input subsets; N, no controlled trial established; NA, survey. \textit{Noise:} S, synthetic intervention; B, recorded-background mixing; F, fixed noisy benchmark; N, no controlled trial established. \textit{Shift:} A, target-adaptive transfer; C, condition transfer with target validation; B, independent-build/layout transfer with limited outcome coverage; T, simulated translation; W, unseen window duration; N, no held-out domain transfer established.


\begingroup\small
\begin{longtable}{@{}P{.25\textwidth}P{.16\textwidth}P{.53\textwidth}@{}}
\caption{Full-Text Appraisal of the 57 Additional Works}\label{tab:expanded-register}\\
\toprule Work & Evidence role & Appraisal focus and qualification \\
\midrule\endfirsthead
\toprule Work & Evidence role & Appraisal focus and qualification \\
\midrule\endhead
P001 \citep{P001} & Context & Single infrared image restoration; no independently
acquired second modality. This is context rather than evidence for a fusion advantage.
\\
P002 \citep{P002} & Context & Infrared machinery diagnosis. Combining thermal regions
remains image-only; vibration fusion is proposed as future work. \\
P003 \citep{P003} & Industrial & Physical camera and audio welding classification. Split
independence and controlled missing/noise evaluation are not established in the
inspected postprint. \\
P004 \citep{P004} & Methods & Seminal evaluation foundation on model-selection bias. It
does not supply an industrial fusion experiment. \\
P005 \citep{P005} & Methods & Transferable modality-dropout method. Training exposure
and planned available-input configurations determine the robustness claim. \\
P006 \citep{P006} & Methods & Transferable gated intermediate fusion. Learned allocation
alone does not establish calibrated modality reliability. \\
P007 \citep{P007} & Methods & Transferable robust multimodal classification
architecture. The inspected preprint and published journal metadata refer to one work.
\\
P008 \citep{P008} & Methods & Predictive uncertainty under dataset shift. Good
in-distribution confidence does not guarantee shifted reliability. \\
P009 \citep{P009} & Methods & Reproducibility and reporting foundation. Documentation
and code availability are distinct from independently reproduced results. \\
P010 \citep{P010} & Methods & The visual patch Transformer was evaluated on image
recognition. The study provides no comparison of the industrial visual models evaluated
here. \\
P011 \citep{P011} & Methods & The audio-spectrogram Transformer represents sound through
spectrogram patches. The present studies include no completed audio-branch comparison.
\\
P012 \citep{P012} & Methods & Self-supervised dilated-convolution time-series encoder.
It is not a Transformer; the reported Transformer substitution performs worse. \\
P013 \citep{P013} & Methods & The video Transformer models interactions across space and
time. The REASSEMBLE visual classifier averages frame descriptors over time. \\
P014 \citep{Nagrani2021MBT} & Methods & Audiovisual attention-bottleneck fusion.
Transferable representation design does not establish industrial sensor-failure
robustness. \\
P015 \citep{P015} & Methods & Multimodal benchmark and evaluation framework. Its
different tasks are not independent replications of one industrial target. \\
P016 \citep{P016} & Methods & The missing-modality Transformer evaluation includes
adverse multimodal results. For the Food dataset, the prose and tabulated unimodal
result disagree. \\
P017 \citep{P017} & Methods & The preprint and published article on leakage in
scientific evaluation give different counts of affected papers. \\
P018 \citep{P018} & Methods & The study evaluates pretrained time-series representations
on benchmarks. Transfer of its weights to the industrial tasks considered here was not
tested. \\
P019 \citep{Sliwowski2025REASSEMBLE} & Dataset & The robotic assembly dataset provides
native action outcomes. Its published demonstration count differs from the eligible
sample for the present action-level analysis. \\
P020 \citep{P020} & Methods & Transferable multimodal negative-learning study. Its noise
regimes and negative comparisons remain specific to the reported tasks. \\
P021 \citep{P021} & Industrial & Physical tool images and force measurements with
adaptive fusion. Post-pass image timing and incompletely reported splits limit
prognostic claims. \\
P022 \citep{P022} & Methods & Non-industrial dynamic-fusion method. Its evidence informs
fusion assumptions without constituting industrial deployment validation. \\
P023 \citep{P023} & Methods & Transferable constrained fusion of text with time series.
It supplies negative evidence against assuming every additional modality helps. \\
P024 \citep{P024} & Industrial & Industrial inspection using geometric prior
information. RGB and geometry are distinguished from independent machine telemetry. \\
P025 \citep{P025} & Industrial & The industrial inspection model combines visual and
geometric features with language for defect reporting. Its evaluation includes
task-specific supervision and measures report consistency. \\
P026 \citep{P026} & Industrial & Train-calibrated RGB/geometry fusion includes explicit
missing-3D evaluation and static/gated controls. It is not camera-plus-machine-telemetry
evidence. \\
P027 \citep{P027} & Industrial & Audio, vibration and current fusion with missing-input
claims. Limited auditability excludes it from strong quantitative synthesis. \\
P028 \citep{Thoduka2024Handover} & Dataset & Completed robotic handover outcomes with
video, force/torque and gripper inputs. Added gripper input reduces one five-class
accuracy result. \\
P029 \citep{Huo2022Multisensor} & Industrial & Thermal-image and vibration fault
diagnosis with single-input and noise controls. Reliability conclusions remain
conditional on its protocol. \\
P030 \citep{Nie2023PatchTST} & Methods & PatchTST was evaluated for time-series
forecasting. The PHM regression and REASSEMBLE classification experiments trained
separate temporal models from scratch. \\
P031 \citep{P031} & Industrial & Vision-language surface localization is evaluated with
pixel metrics. These measure localization rather than image-level detection or numerical
process sensing. \\
P032 \citep{Hua2026Weld} & Industrial & Camera/audio weld-penetration attention fusion.
Clean single-input and fused models all report 100\% accuracy; independent-run splitting
is unclear. \\
P033 \citep{P033} & Industrial & RGB/geometry disagreement scoring. Eyecandies image
AUROC decreases despite improved localization, showing metric-specific trade-offs. \\
P034 \citep{P034} & Methods & Space-time attention models cross-frame interactions. The
core REASSEMBLE visual representation averages frame descriptors. \\
P035 \citep{P035} & Context & Weakly supervised single-camera surface inspection.
Internal feature fusion is not fusion of independently acquired physical modalities. \\
P036 \citep{P036} & Methods & Transferable counterargument on attention interpretation.
Model inspection does not by itself establish a physical causal explanation. \\
P037 \citep{Altinses2023Fuzzy} & Simulation & Camera/robot telemetry corruption in
simulated pose regression. Overlapping temporal windows and random splitting limit
physical-deployment inference. \\
P038 \citep{P038} & Industrial & RGB/geometry consistency fusion. Its learned scalar
mixing coefficient does not demonstrate sample-conditioned online reliability
adaptation. \\
P040 \citep{Bengio2004CVVariance} & Methods & Cross-validation variance theory.
Overlapping training folds do not supply independent replications for a naive standard
error. \\
P041 \citep{Guo2017Calibration} & Methods & Validation-fitted calibration foundation.
In-distribution temperature scaling is not a guarantee of reliable probabilities after
shift. \\
P042 \citep{P042} & Context & Single-vibration diagnosis with synthetic noise. Simulated
target training and reporting inconsistencies limit transfer and field-validity claims.
\\
P043 \citep{P043} & Methods & Distribution-shift benchmark. OOD validation access is
explicit; this is not an untouched-target industrial multimodal experiment. \\
P044 \citep{P044} & Methods & Transferable medical fusion counterexample. Selection and
evaluation share validation data; no independent final test is established. \\
P045 \citep{P045} & Industrial & Post-pass tool photographs and acoustic emission
predict wear. Strong visual and adverse multimodal comparators coexist with unresolved
tool-split coverage. \\
P046 \citep{P046} & Industrial & Microphone/accelerometer bearing fusion. Test-noise
results lack an equivalent static corruption comparator; recording-disjoint splitting is
unclear. \\
P048 \citep{P048} & Industrial & Thermal-camera and vibration feature fusion.
Overlapping windows, multirate pairing and separately fitted speeds restrict
generalization claims. \\
P049 \citep{P049} & Methods & Transferable attention-interpretability critique.
Attention perturbation is distinct from physical input corruption or modality dropout.
\\
P050 \citep{P050} & Methods & Transferable unaligned language/audio/video
cross-attention. Unequal feature rates are not a controlled hardware-clock offset test.
\\
P051 \citep{P051} & Industrial & Visible/infrared federated detection. Incomplete
clients differ from surprise inference outage; group splits, clipped results and
coefficient conflicts limit appraisal. \\
P052 \citep{Kim2026MSM} & Dataset & Milling dataset with laboratory video and machine
measurements. Dataset availability enables research but is not itself a learned fusion
result. \\
P053 \citep{P053} & Industrial & Audio/vibration wavelet-image cross-attention. These
are transformed signals, not camera images; controlled missing/noise tests are not
established. \\
P054 \citep{P054} & Industrial & RGB/surface-normal anomaly detection. Training
feature-noise regularization is distinct from a test-input degradation benchmark. \\
P055 \citep{P055} & Methods & Proper-scoring-rule foundation. Expected propriety does
not certify finite-sample calibration or shifted sensor reliability. \\
P056 \citep{P056} & Industrial & RGB/point-cloud local alignment preprint, version 1.
Checkpoint-selection population and retraining variability remain unclear. \\
P057 \citep{Zeng2026ViFailback} & Dataset & The study uses visual symbols and language
to diagnose manipulation failures. The present ViFailback experiment instead combines
camera views and numerical robot observations. \\
P058 \citep{Zhang2023QMF} & Methods & Transferable quality-aware weighting. Its
conditional generalization-upper-bound comparison does not imply universal actual-risk
superiority. \\
P059 \citep{Zhao2024RTDETR} & Methods & Single-image object detector foundation.
Executing an extracted CNN backbone does not execute the full detector or its reported
inference pipeline. \\
\bottomrule\end{longtable}
\displaynote{The accompanying CSV identifies each target, fusion method, control, test
of missing, noisy, or shifted inputs, result, and source location. The role column
distinguishes direct experimental evidence from background material.}\endgroup
\noindent\textit{Interpretation:} Transferable methods are non-industrial unless their
record identifies a direct application. Dataset papers and unimodal/context references
do not supply a multimodal fusion experiment. A retained preprint remains provisional.
All source versions, complete qualifications and numerical locators are in the expanded
register.


\section{Benchmark Identity and Assessment Populations}
The MVTec 3D-AD dataset paper lists 249 anomaly-free and 948 anomalous test samples, a total of 1,197. It reports 2,656 training and 294 validation samples separately \citep[Table~1, PDF p.~4]{Bergmann2022The}. UCF reports the same test count \citep{Cheng2025Multimodal}, whereas some other method descriptions give 1,137. The discrepancy leaves the samples underlying those comparisons uncertain. The review uses the published scores, whose comparability also depends on preprocessing, masks and metric definitions.
